\section{Cancellation of Gamma factors before finite expansion}
\label{app:gamma-cancel}
The connected reduction in Appendix~\ref{app:loop-tree} expresses the coefficient as a ratio of products of Gamma functions. Every argument has the form $A+Bd/2$, where $A$ and $B$ are integers. The tree fixes $B$, while the target powers determine $A$. We first explain how to reuse this fixed part and then derive the cancellation rule.

Let $n_j$ be the number of vertices in the subtree rooted at a nonroot vertex $j$, including $j$ itself. Write the integer part of its incoming effective power as $U_j$. The definitions in Eq.~\eqref{eq:loop-tree-transfer} give
\begin{equation}
 u_j=U_j-(n_j-1)\frac d2,
 \qquad U_j=a_j+\sum_{h\in\mathcal C(j)}(U_h+v_h),
 \qquad \sigma_j=U_j+v_j-n_j\frac d2.
 \label{eq:gamma-cancel-subtree}
\end{equation}
For a leaf, $n_j=1$ and $U_j=a_j$, so this expression agrees with the one-bubble formula. If it holds for the children, summing their transferred powers gives the expression at their parent because $n_j=1+\sum_h n_h$. Induction therefore proves the formula throughout the tree. The subtree sizes can be calculated once for a batch of targets. At the root, define $U_0=a_0+\sum_{h\in\mathcal C(0)}(U_h+v_h)$, so the corresponding power is $\alpha=U_0-(L-1)d/2$.

Substitution in the bubble factor gives numerator arguments
\begin{equation}
 U_j+v_j-n_jd/2,\qquad -U_j+n_jd/2,\qquad -v_j+d/2,
\end{equation}
and denominator arguments
\begin{equation}
 U_j+(1-n_j)d/2,\qquad v_j,\qquad -U_j-v_j+(n_j+1)d/2.
\end{equation}
The root-triangle arguments follow by substituting $\alpha$ into Eq.~\eqref{eq:loop-tree-H}. These expressions evaluate the same Gamma arguments as the preceding recursive method and avoid calculating their dimension-dependent parts separately for each target.

After dividing by the unit-power master, collect the factors with a fixed value of $B$. Cancel identical arguments and order the remaining integer parts as $A_1<\cdots<A_m$. Let $w_i$ be the signed multiplicity of $\Gamma(A_i+Bd/2)$, with numerator factors positive. The target and master have matching argument slopes, so $\sum_i w_i=0$. Define $C_i=\sum_{j=1}^i w_j$. Then
\begin{equation}
 \prod_{i=1}^{m}\Gamma(A_i+Bd/2)^{w_i}
 =\prod_{i=1}^{m-1}\prod_{k=A_i}^{A_{i+1}-1}
       (k+Bd/2)^{-C_i}.
 \label{eq:gamma-cancel-product}
\end{equation}
To prove this identity, use Eq.~\eqref{eq:loop-tree-shift} to write every Gamma function relative to $\Gamma(A_1+Bd/2)$. Its total exponent is zero. A factor in the interval from $A_i$ to $A_{i+1}-1$ appears in precisely the terms with index greater than $i$, so its exponent is $\sum_{j>i}w_j=-C_i$. This gives Eq.~\eqref{eq:gamma-cancel-product}. Intervals with $C_i=0$ contribute one and require no factors. The identity holds meromorphically under the same physical assumptions as the connected reduction.

Finally, we normalize each linear factor by the integer greatest common divisor of its two coefficients and combine equal normalized factors. Integer numerator and denominator products accumulate separately, and one rational simplification reduces their ratio. We compute the reference-master arguments once per target batch. These changes preserve the exact coefficient and master basis.

The reserved comparison uses the distinct star, chain, and branched structures defined in Appendix~\ref{app:loop-tree}, giving 34 structures through thirteen loops. For each structure, the small batch uses the seven stated target assignments. The large batch includes those assignments and additional distinct powers: every positive physical power lies between one and four, and the middle triangle power lies between minus one and four. Auxiliary powers remain zero. The scale $s$ stays symbolic. Each large batch has 128 targets except at one loop, where the allowed set contains 96 targets. The selected finite sets are recorded with the results; the analytic formula applies throughout the stated domain.

Five fresh processes measure each method on each request. The comparison rotates method order and uses the same input, output procedure, machine, and master basis. The ratio divides sums of request medians, with request-resampling intervals defined in Eq.~\eqref{eq:loop-tree-resampling}. Loading, recognition, coefficient calculation, output, and process exit enter the timer. Exact comparison with the preceding coefficient routine and independent Gamma evaluations occur after timing and are excluded for both methods. The four regression requests exercise two earlier analytic families at thirteen loops and planar and nonplanar two-loop native batches. They test the unchanged fallback path and do not provide a new estimate of the separate ML contribution.
